\documentclass[../main.tex]{subfiles}
\begin{document}

This appendix establishes the two things the main experiment of Chapter~\ref{chap:results} builds on. The first is the evaluation protocol used throughout that chapter, the single selection metric, the development dataset, and the training loss, each fixed by a baseline experiment rather than assumed. The second is the evidence that motivates the proposed hybrid model: configurations with near-identical scores disagree on which clips they classify, and the OpenFace and I3D feature families fail on different clips, so learning to fuse them can exceed any single encoder. The baseline rankings themselves enter only as the reference level Chapter~\ref{chap:results} must exceed.

\FloatBarrier
\section{Experimental Setup}
\label{sec:setup-baseline}

Every result in this appendix comes from one apparatus. The five single-encoder baselines of Section~\ref{sec:baselines} are each trained under the three losses of Section~\ref{sec:losses-method} on each of the three training sources, optimised identically (Adam, fixed seed, early stopping on validation loss; Section~\ref{sec:training}) and evaluated on the held-out test split under the 3$\times$3 cross-dataset protocol of Section~\ref{sec:eval}. The in-domain experiments (Sections~\ref{sec:protocol}--\ref{sec:agreement}) use the fifteen CMOSE-trained, CMOSE-tested configurations (five architectures $\times$ three losses); the cross-dataset experiment (Section~\ref{sec:crossdomain}) uses the full matrix.

Each run is scored on all six metrics of Section~\ref{sec:metrics-lit}, but the reporting is deliberately reduced to follow the protocol. The first experiment (Section~\ref{sec:protocol-metrics}) settles on QWK as the single metric that selects every best model; from then on the two balanced secondary metrics, macro-accuracy and macro-MAE, accompany each headline result to describe how a model treats the rare classes, while accuracy, micro-MAE, and Cohen's kappa are reported only where they are themselves the object of study, accuracy where its unreliability is the point (Section~\ref{sec:closing-loop}) and Cohen's kappa where prediction agreement is (Section~\ref{sec:agreement}). The overview figures that survey a full design space or the whole cross-dataset matrix (Figure~\ref{fig:base_models_all_metrics} and Figure~\ref{fig:crossdomain_base} here, Figure~\ref{fig:hybrid_ablation_all_metrics} in Chapter~\ref{chap:results}) show the complete six-metric grid for reference; their discussion, and every other result, follows this reduced reporting.

Two aggregation conventions hold across this appendix and the experimental results (Chapter~\ref{chap:results}). A winner is the single configuration with the highest QWK over the tunable axis (here the loss), ties broken stable-first: because QWK is the only selection metric, a reported best model never depends on any other number. A claim about a family of configurations is instead reported as a mean and a distribution over the design grid. Every configuration is trained once with a fixed seed (Section~\ref{sec:training}), so a reported mean measures robustness to design choices, not noise across seeds.

\FloatBarrier
\section{Evaluation Protocol}
\label{sec:protocol}

Before the baselines can be ranked, three parameters of the protocol must be fixed: the metric that selects models, the dataset the models are developed on, and the loss they are trained with. Each is settled by one experiment below. Two of the arguments need out-of-domain evidence and are stated here as claims, then proved in Section~\ref{sec:closing-loop}.

\subsection{Selecting the Primary Metric}
\label{sec:protocol-metrics}

The first parameter to fix is the metric that selects models, because on an ordinal, imbalanced task that choice is not neutral. Engagement levels lie on a line, and 69.4\% of CMOSE is the single class E2 (Engage) (Table~\ref{tab:T3_1_dataset_stats}); accuracy ignores both properties, rewarding majority-class collapse and counting a one-level miss the same as a three-level miss, whereas QWK is chance-corrected and weights every error by the squared distance between predicted and true levels (Section~\ref{sec:metrics-lit}). To test the consequence, the fifteen in-domain baseline configurations are scored on all six metrics, and two questions are put to them: does the leaderboard winner depend on the metric, and do the metrics rank the configurations the same way.

The winner depends on the metric. Table~\ref{tab:T4_1_per_metric_winner} reports which of the fifteen configurations each metric declares best in-domain, and no two camps agree: QWK picks the OpenFace TCN under cross-entropy (0.537), both macro metrics pick the I3D MLP under the ordinal loss (macro-accuracy 0.635, macro-MAE 0.458), and accuracy, micro-MAE, and Cohen's kappa all pick the I3D MLP under cross-entropy (0.773, 0.244, 0.436).
\begin{table}[ht]
\centering
\caption{Winning baseline configuration by metric (in-domain CMOSE).}
\label{tab:T4_1_per_metric_winner}
\fittable{%
\begin{tabular}{lllr}
\toprule
Metric & Best model & Loss & Value \\
\midrule
QWK & openface\_tcn & CE & 0.537 \\
Macro-Accuracy & i3d\_mlp & Ordinal & 0.635 \\
Macro-MAE & i3d\_mlp & Ordinal & 0.458 \\
Accuracy & i3d\_mlp & CE & 0.773 \\
MAE & i3d\_mlp & CE & 0.244 \\
Cohen $\kappa$ & i3d\_mlp & CE & 0.436 \\
\bottomrule
\end{tabular}%
}
\end{table}


The metrics also fall into two blocks bridged only by QWK. Figure~\ref{fig:metric_correlation_base} reports the Kendall rank correlation between every pair of metrics over the fifteen configurations: a micro block (accuracy and micro-MAE, $\tau=0.98$, dominated by the majority class) and a macro block (macro-accuracy and macro-MAE, $\tau=0.89$) that are almost unrelated to each other (cross-block $\tau\approx0.01$--$0.10$). QWK is the only metric with moderate-to-strong agreement with both blocks ($\tau=0.49$--$0.60$).

\begin{figure}[ht]
\centering
\includegraphics[width=0.62\linewidth]{metric_correlation_base.png}
\caption{Rank agreement (Kendall $\tau$) between the six metrics over the fifteen
baseline configurations, in-domain on CMOSE (MAE metrics sign-flipped so positive $\tau$
means~agreement).}
\label{fig:metric_correlation_base}
\end{figure}

The choice of metric therefore changes the result and must be argued, not defaulted. This thesis adopts QWK as its single primary metric, the compromise sensitive to ordinal distance without being dominated by the majority class: the winner is the argmax of QWK, ties broken stable-first, with no second metric consulted. The remaining five are secondary. The decisive argument against accuracy, that it cannot even detect the collapse of a model's ordinal signal under shift, needs out-of-domain evidence and is deferred to Section~\ref{sec:closing-loop}. With the metric fixed, the next question is which dataset to develop on.

\subsection{Selecting the Development Dataset}
\label{sec:cmose-anchor}

A development dataset is usable only if models rise meaningfully above chance on it, since a dataset where they do not cannot separate good architectures from bad, and only CMOSE clears that bar. The best in-domain model reaches QWK 0.537 on CMOSE but only 0.166 on DAiSEE, while DAiSEE's raw accuracy of 0.548 hides the failure, a first sign of how misleading accuracy can be. CMOSE is therefore the development dataset; DAiSEE re-enters as a transfer source and target in Section~\ref{sec:crossdomain} and in Chapter~\ref{chap:results}, where its difficulty is informative rather than obstructive. With the metric and the dataset fixed, the last parameter is the loss.

\subsection{Selecting the Development Loss}
\label{sec:loss-tradeoff}

The loss is the axis that trades the primary metric against the balanced secondary ones, so it must be fixed before the architecture is judged. For every baseline, the loss is swept from cross-entropy through weighted cross-entropy to the ordinal loss, in-domain on CMOSE, and QWK is plotted against the balanced metrics (Figure~\ref{fig:loss_metric_tradeoff}).

\begin{figure}[ht]
\centering
\includegraphics[width=\linewidth]{loss_metric_tradeoff.png}
\caption{QWK against the balanced secondary metrics (macro-accuracy and macro-MAE) for
every baseline across the three losses, in-domain on CMOSE.}
\label{fig:loss_metric_tradeoff}
\end{figure}

The trade is Pareto: rebalancing the loss raises macro-accuracy and improves macro-MAE while QWK falls for four of the five baselines, so no loss is best on QWK and on the balanced metrics at once. For the I3D MLP, switching from cross-entropy to the ordinal loss lifts macro-accuracy from 0.526 to 0.635 and improves macro-MAE from 0.585 to 0.458 (the best balanced figures of any baseline) while QWK falls from 0.519 to 0.487. Cross-entropy is the development loss because it maximises QWK, the selection metric, and isolates the effect of the architecture from that of the loss (best-architecture QWK by loss: CE 0.537, weighted CE 0.500, ordinal 0.487); the proposed hybrid model of Chapter~\ref{chap:results} is consequently trained with it. That this CE $>$ weighted CE $>$ ordinal ranking survives domain shift is the second deferred claim, closed in Section~\ref{sec:closing-loop}. With the three parameters fixed, the baselines can be ranked.

\FloatBarrier
\section{In-Domain Results}
\label{sec:baseline-results}

The aim of this experiment is to set the reference level the proposed architecture must exceed. The five single-encoder baselines are trained and tested in-domain on CMOSE under each of the three losses (four on the OpenFace descriptor: an MLP, a TCN, an LSTM, and a Transformer; and one on the I3D clip vector, the I3D MLP); Figure~\ref{fig:base_models_all_metrics} reports the complete grid on all six metrics for reference.

\begin{figure}[ht]
\centering
\includegraphics[width=\linewidth]{base_models_all_metrics.png}
\caption{The five baseline models under the three losses on all six metrics, in-domain~on~CMOSE.}
\label{fig:base_models_all_metrics}
\end{figure}

On QWK, the selection metric, the strongest baseline is the OpenFace TCN under cross-entropy at QWK 0.537: this is the reference level. The grid adds two structural observations. The loss axis determines on which metric family a model is strongest, cross-entropy on the micro and agreement metrics, the rebalanced losses on the macro metrics that weight the rare disengaged classes equally. The architecture axis determines which kind of feature is strongest: the OpenFace TCN's per-frame facial dynamics (a head turning away, gaze drifting off-screen) separate the rarer disengaged levels and rank highest on QWK, while the I3D MLP, over a single pooled appearance vector, aligns with the majority E2 class and ranks highest on every micro metric. On QWK the two score almost the same (0.537 versus 0.519). Two models this close, yet built on disjoint feature families, pose a question the score table cannot answer: are they interchangeable, or do they succeed on different clips?

\FloatBarrier
\section{Prediction Agreement}
\label{sec:agreement}

Scores summarise how often a model is right, not which clips it is right on. This experiment descends to individual predictions on the 1221 in-domain CMOSE test clips to ask whether models with similar scores make the same decisions, first across the whole grid, then for the two feature families. Cohen's kappa, the agreement statistic, is the metric of study here.

The fifteen configurations agree far less than their clustered scores suggest. Figure~\ref{fig:agreement_base_models} reports the pairwise Cohen's kappa between every pair of predictions: mean pairwise agreement is only $\kappa=0.27$--$0.39$, and sharing the architecture ($\kappa=0.391$) or the loss ($\kappa=0.374$) raises it only mildly over sharing neither ($\kappa=0.273$). The configurations are different hypotheses about the data, not reparametrisations of one another.

\begin{figure}[ht]
\centering
\includegraphics[width=\linewidth]{agreement_base_models.png}
\caption{Pairwise prediction agreement (Cohen's $\kappa$) between all fifteen baseline
configurations on the in-domain CMOSE test set, in model-major blocks.}
\label{fig:agreement_base_models}
\end{figure}

The consequence is headroom. The best single configuration classifies 77.3\% of the test clips correctly, but at least one of the fifteen is correct on 97.6\% (the oracle), so only 2.4\% of clips defeat every model. Plain majority voting over the fifteen lands at 73.5\%, below the best single model, because the errors are too decorrelated for unweighted voting to use. Claiming the headroom therefore requires a learned combination of complementary views of the input, which is the design premise of the proposed hybrid model of Chapter~\ref{chap:results}: it fuses separately encoded behavioural streams instead of averaging finished predictions.

The two feature families scored almost the same, yet differ in nature: OpenFace is a per-frame, structured account of facial behaviour (gaze, eye and face landmarks, head pose, action units), while I3D is a single clip-level vector of global appearance and motion. Whether their close scores mean redundancy or complementarity is the question, and by scores alone the I3D MLP appears merely redundant, since the OpenFace TCN edges it on QWK and on both balanced secondary metrics. The agreement map says otherwise: in Figure~\ref{fig:agreement_base_models} the I3D MLP's row is the palest in the grid, its agreement with the OpenFace TCN among the lowest of any pair despite the two having the closest scores. The clip-level decomposition confirms it: against the OpenFace TCN, the I3D MLP has the largest exclusive-correct share of any pair in the grid (8.6\% of test clips only I3D gets right) and the smallest both-wrong share (15.3\%), and the pair's oracle accuracy is 84.7\%, $+7.4$ points over the best single model and the largest pair gain in the grid. The behaviour descriptors and the appearance features look at different evidence and fail on different clips, which is the empirical motivation for fusing the two streams in the proposed hybrid model of Chapter~\ref{chap:results} that a score table alone could never supply.

\FloatBarrier
\section{Cross-Dataset Generalisation}
\label{sec:crossdomain}

In-domain numbers are not the deployment question, and two of the protocol arguments were deferred to out-of-domain evidence. Training each baseline on one dataset and testing on another produces the baseline cross-dataset matrix that Chapter~\ref{chap:results} measures the proposed model against, and it closes those two deferred arguments: whether the metric and the loss fixed in-domain remain the right choices once the ordinal signal collapses under a change of dataset.

\begin{figure}[ht]
\centering
\includegraphics[width=\linewidth]{crossdomain_base.png}
\caption{Cross-domain generalisation of the best baseline per cell, on all six metrics; each
cell is the single QWK-selected best baseline, reported on every metric.}
\label{fig:crossdomain_base}
\end{figure}

The 3$\times$3 heatmaps of Figure~\ref{fig:crossdomain_base} are the baseline that Chapter~\ref{chap:results} measures the proposed model against, and they show a cross-dataset collapse. On the diagonal QWK is healthy (CMOSE 0.54), macro-accuracy is above chance (CMOSE 0.54), and ordinal error is moderate (CMOSE macro-MAE 0.55). Off the diagonal all three collapse together: CMOSE$\rightarrow$DAiSEE QWK $=0.02$ and DAiSEE$\rightarrow$CMOSE $=0.05$ (chance-level agreement), macro-accuracy falls toward the 0.25 chance floor, and macro-MAE inflates past a full class. The collapse is roughly symmetric, which rules out one dataset being a superset of the other and points to a real shift in both label priors and recording conditions. This is the regime in which the two protocol arguments deferred above can finally be tested.

\phantomsection\label{sec:closing-loop}
The cross-dataset matrix supplies the out-of-domain evidence the protocol deferred. Accuracy fails to register the collapse, which completes Section~\ref{sec:protocol-metrics}. On DAiSEE$\rightarrow$CMOSE the QWK reads 0.05, chance-level ordinal agreement, yet a model that simply predicted CMOSE's 69\% majority class (Table~\ref{tab:T3_1_dataset_stats}) would post 0.69 accuracy, higher than the QWK-selected baseline's 0.35 in that cell (accuracy panel of Figure~\ref{fig:crossdomain_base}). Accuracy would therefore rank a majority-collapsed predictor above the model that actually tracks the ordinal structure, staying comfortable precisely where the ordinal signal is dead. That makes accuracy the least reliable metric here and QWK, the only bridge between the micro and macro blocks, the most.

The loss ranking is preserved under the shift, which completes Section~\ref{sec:loss-tradeoff}. On the unseen-target cells the best-architecture QWK by loss follows the same CE $>$ weighted CE $>$ ordinal order as in-domain. The loss ranking transfers where the architecture ranking does not, so choosing cross-entropy in-domain remains right out of domain, confirming after the fact the choice to establish it as the development~loss.

\end{document}
